\begin{tikzpicture}[
  every node/.append style={font=\figurefont\footnotesize},
  source/.style={book node,fill=FigureInput!12,draw=FigureInput,
    text width=38mm,minimum height=42mm,align=left},
  task/.style={book node,fill=FigureModel!10,draw=FigureModel,
    text width=23mm,minimum height=12mm,align=center},
  result/.style={book node,fill=FigureOutput!10,draw=FigureOutput,
    text width=65mm,minimum height=12mm,align=left}
]
  \node[source] (notice) at (21mm,0) {
    \textbf{教学通知}\\
    ML-0418\\
    2025年4月15日发布\\
    读书会改至4月18日\\
    14:00，海棠楼201。\\
    校外来宾须在\\
    4月16日18:00前报名。
  };
  \node[task] (analysis) at (67mm,54mm) {分析结构};
  \node[task] (classify) at (67mm,36mm) {类别};
  \node[task] (extract) at (67mm,18mm) {事实记录};
  \node[task] (summary) at (67mm,0) {摘要};
  \node[task] (translation) at (67mm,-18mm) {译文};
  \node[task] (answer) at (67mm,-36mm) {答案};
  \node[task] (speech) at (67mm,-54mm) {朗读};

  \node[result] (analysis-r) at (121mm,54mm)
    {“须报名”的义务对象是校外来宾};
  \node[result] (classify-r) at (121mm,36mm)
    {活动通知、时间变更};
  \node[result] (extract-r) at (121mm,18mm)
    {新时间：2025-04-18 14:00\\地点：海棠楼201};
  \node[result] (summary-r) at (121mm,0)
    {保留新时间、地点和精确报名条件};
  \node[result] (translation-r) at (121mm,-18mm)
    {Friday, April 18, at 14:00\\Room 201, Haitang Building};
  \node[result] (answer-r) at (121mm,-36mm)
    {“明天能来吗？”需日期、身份及报名状态};
  \node[result] (speech-r) at (121mm,-54mm)
    {“四月十八日，下午两点”\\输出为可播放声音，不能只交付文字};

  \draw[NNLine,line width=0.85pt] (notice.east) -- (48mm,0);
  \draw[NNLine,line width=0.85pt] (48mm,-54mm) -- (48mm,54mm);
  \foreach \name in {analysis,classify,extract,summary,translation,answer,speech} {
    \draw[book arrow] (48mm,0 |- \name.west) -- (\name.west);
    \draw[book arrow] (\name.east) -- (\name-r.west);
  }
\end{tikzpicture}
